%%%
%
% $Autor: Shubhankar Dumka, Gautam Ramesh, Siddhanth Sharma $
% $Datum: 2025-10-06 $
% $Dateiname: 
% $Version: 1 $
%
% !TeX spellcheck = GB
% !TeX program = pdflatex
% !TeX encoding = utf8
%
%%%

\chapter{\textcolor{red}{Description of Important Modules}}

\index{Modules}

% --------------------------------------------------------
\noindent
Computer vision systems are reliant on efficient integration and implementation of libraries that are able to handle image processing tasks seamlessly, along with deep learning frameworks and algorithms to achieve accuracy in performance. From the context of this project, \textbf{OpenCV}, \textbf{PyTorch}, and \textbf{Convolutional Neural Networks (CNNs)} form the foundation for implementing license plate detection systems on the NVIDIA Jetson Nano B01. Combined, such components allow for the creation of a complete end-to-end pipeline that can handle tasks relevant to image acquisition, pre-processing and deep learning inference.

The \textbf{OpenCV (Open Source Computer Vision Library)} functions as core image and video processing framework. OpenCV is an open-source library that offers a diverse range of algorithms that are optimized real-time computer vision tasks \cite{Bradski:2000,opencv_wiki}. It covers tools for camera interfacing, frame capture, color space conversion, resizing, and visual annotation of detection results. Due to it being optimized for hardware acceleration, OpenCV is specifically suited for embedded platforms where computational resources are constrained but real-time image processing is required \cite{opencv_lfs,opencv_on_wheels}.

In conjunction to this image-processing layer, another important library is \textbf{PyTorch}, a modern and open-source deep learning framework. PyTorch works on a dynamic, define-by-run execution model that enables neural networks construction, modification during runtime and facilitates efficient debugging \cite{paszke2019pytorch}. At the core of its definition, PyTorch enables usage of a powerful tensor library with GPU acceleration and automatic differentiation via its \texttt{autograd} engine. Deep learning workloads are simplified through efficient gradient computation \cite{pytorchdocs}. PyTorch enables the translation of pre-processed image frames into tensors, leveraging the trained detection model on the Jetson Nano’s CUDA-enabled GPU, and inference workflow management through libraries such as \texttt{torchvision} \cite{jetsonpytorch}. 

At the algorithmic core of the detection pipeline are \textbf{Convolutional Neural Networks (CNNs)}, forming the basis for implementing deep learning architectures pertaining to image-based tasks. CNNs are essentially designed to compute data with grid-like structures, like images, by exploiting local spatial correlations through convolutional layers and parameter sharing \cite{lecun1998gradient,goodfellow2016deep}. Characteristics such as edges, textures and complex semantic patterns are extracted in a sequential manners by CNNs. This aids in object detection through varying lighting environments and viewing angles \cite{krizhevsky2012imagenet,simonyan2014very}. 

Overall, OpenCV, PyTorch, and CNNs form a reliable processing pipeline: OpenCV takes care of image acquisition and preprocessing, PyTorch provides the computational framework for model execution, and CNNs provide the representational power suitable for reliable visual detection. Such a modular and coherent architecture enables real-time computer vision on embedded hardware and offers a foundation for applications in traffic monitoring, security, and intelligent transportation systems \cite{jetsonpytorch,nvidia_jetson}. Subsequent chapters describe these components comprehensively, offering details to their structure, installation, and role within the overall system.
